%%%%%%%%%%%%%%%%%%%%%%%%%%%%% Define Article %%%%%%%%%%%%%%%%%%%%%%%%%%%%%%%%%%
\documentclass{article}
%%%%%%%%%%%%%%%%%%%%%%%%%%%%%%%%%%%%%%%%%%%%%%%%%%%%%%%%%%%%%%%%%%%%%%%%%%%%%%%

%%%%%%%%%%%%%%%%%%%%%%%%%%%%% Using Packages %%%%%%%%%%%%%%%%%%%%%%%%%%%%%%%%%%
\usepackage{listings, xcolor}
\usepackage{parskip}
\usepackage{hyperref}
\usepackage{amsmath}
\usepackage{graphicx}
\usepackage[a4paper, margin = 1.5cm]{geometry}
\usepackage{float}
\usepackage{amssymb}

%%%%%%%%%%%%%%%%%%%%%%%%%%%%%%% Title & Author %%%%%%%%%%%%%%%%%%%%%%%%%%%%%%%%
\title{Diagramma delle Classi - Costruzione}
\author{Trentin Tommaso}
%%%%%%%%%%%%%%%%%%%%%%%%%%%%%%%%%%%%%%%%%%%%%%%%%%%%%%%%%%%%%%%%%%%%%%%%%%%%%%%

\begin{document}
  \section{Livelli di dettaglio}
    Diagrammi UML possono essere creati in diverse fasi del ciclo di vita, con diversi scopi / livelli di dettaglio:
    \begin{itemize}
      \item \textbf{System domain model}: modella entità del dominio del problema che saranno presenti nel sistema (dati di interesse), può contenere le responsabilità delle classi (opzionale); 
      \item \textbf{System model}: include anche classi che saranno usate per costruire il sistema completo (classi per interfaccia utente, classi associate a parti di architettura, classi di utility), oltre a metodi e informazioni sulla navigabilità. 
    \end{itemize}
  \section{Costruzione del System Domain Model}
    \begin{enumerate}
      \item Identificare primo insieme di \textbf{classi} candidate; 
      \item Aggiungere \textbf{associazioni} ed \textbf{attribuuti} delle classi;
      \item Trovare \textbf{generalizzazioni};
      \item Trovare principali \textbf{responsabilità} di ogni classe;
      \item Iterare processo finché modello ottenuto è soddisfacente.
    \end{enumerate}
    \subsection{Identificare Classi}
      Classi dovrebbero corrispondere a entità del dominio del problema, si applicano i concetti di OO: una classe identifica un tipo di dato astratto. Una buona classe:
      \begin{itemize}
        \item ha un nome che ne rispecchia l'intento; 
        \item è un'astrazione che modella un elemento del dominio del problema;
        \item ha insieme ridotto e ben definito di responsabilità.
      \end{itemize}
      \subsubsection{Analisi dei nomi}
        Semplice tecnica per scoprire le classi del dominio:
        \begin{enumerate}
          \item Analizzare documentazione di partenza; 
          \item Elencare i nomi; 
          \item Eliminare nomi che:
            \begin{itemize}
              \item sono ridondanti; 
              \item rappresentano / caratterizzano istanze e non classi;
              \item sono vaghi / generici;
              \item corrispondono a classi non necessarie al livello considerato.
            \end{itemize}
        \end{enumerate}
    \subsection{Identificare Attributi}
    \subsection{}
    \subsection{}
    \subsection{}
\end{document}
